\textbf{What remains for RQ1.}
The pipeline, the benchmark construction and all six methods are implemented and run end to end (Section~\ref{sec:experiment}). What remains is to run the graph part of the benchmark on the real trees:
\begin{enumerate}
\item \emph{Relational export.} Export the family links of all trees into the same versioned Parquet snapshot as the person records (\texttt{acadian-er snapshot}, using the library loaders). The export code is written and tested on synthetic data in the same format.
\item \emph{Full benchmark on real data.} Rebuild the RQ1 benchmark from the full snapshot with all six tiers (t3--t5 need relatives), with anchors that have at least three relatives in their tree, and train the R-GCN at $k=0,1,2$ on the real graph. This also restores the relatives-based criteria of the \texttt{exact\_name} baseline, so that it is compared at full strength.
\item \emph{Clone depth sensitivity.} Cloning only the immediate family makes a clone's 2-hop neighbourhood thinner than the original's (no grandparents, siblings or in-laws), which works against $k=2$. We will report results with clone depth~1 (main result) and depth~2 side by side.
\item \emph{Analysis.} Report per-tier results with bootstrap confidence intervals over queries, and inspect failure cases per tier.
\end{enumerate}
No hyperparameter tuning has been done; we will tune only on a validation split of anchors, disjoint from the test anchors.

\textbf{RQ2.}
The verification pipeline (fact-identified contexts at $k=0$ and $k=2$, prompt, JSON parsing, faithfulness check, metrics, and local and API LLM backends) is implemented and tested with a deterministic fake model. Once the relational export is done we will draw candidate pairs from the RQ1 rankers' top results on real data, hand-label 150--200 of them (\textit{match}, \textit{not match}, \textit{unsure}), and run the LLMs at temperature~0 (the choice of models is still open: a local model through Ollama and/or a hosted one through an OpenAI-compatible API).

\textbf{Changes from the proposal, and why.}
\begin{itemize}
\item \emph{Two experiments, not a pipeline.} As clarified with the instructor after Milestone~1, RQ1 and RQ2 are two separate experiments with the same $k=0$ vs.\ $k>0$ contrast; RQ1 only supplies RQ2's candidate pairs.
\item \emph{Controlled noise.} RQ1 robustness is measured with explicit perturbation tiers rather than relying on noise in the data, so each result can be attributed to one kind of noise.
\item \emph{An attribute-only ablation of our own model.} We added the R-GCN at $k=0$ (same model, features and training, no neighbours), so that any gain at $k=1,2$ can be attributed to the neighbourhood and not to the training procedure. The baselines are now the existing criteria method, a TF-IDF and a sentence-transformer encoder.
\item \emph{Faithfulness instead of RAGAS.} RAGAS scores free-text answers with another LLM. Because every fact in our contexts has an identifier, we can check citations exactly (does the cited fact exist, with the cited value?), which is cheaper, deterministic and easier to interpret.
\item \emph{Related work.} We replaced the GraphRAG summarisation and spatiotemporal imputation papers with work closer to our two RQs: structural linkage of historical records, the R-GCN itself, and LLM-based entity matching.
\end{itemize}
